% !TEX program = xelatex
% ============================================================================
%  POSN Computer Qualification — Logic Slide Deck
%  Topic: ตรรกศาสตร์พื้นฐาน (Basic Logic)
% ============================================================================

\documentclass[11pt, aspectratio=169]{beamer}

% --- Load shared preamble ---
% ============================================================================
%  POSN Computer Qualification — Shared Beamer Preamble
%  Used by all topic slide decks. Compile with XeLaTeX.
% ============================================================================

% ---------- Thai language & font ----------
\usepackage{iftex}
\usepackage{fontspec}

% XeTeX-specific: enable Thai line breaking
\ifXeTeX
  \XeTeXlinebreaklocale{th}
\fi
% Noto Sans Thai — body text (clean modern sans-serif, handles all Thai)
\setmainfont{Noto Sans Thai}[
  UprightFont = Noto Sans Thai Light,
  BoldFont    = Noto Sans Thai Medium,
  Scale       = 0.9
]
\setsansfont{Noto Sans Thai}[
  UprightFont = Noto Sans Thai Light,
  BoldFont    = Noto Sans Thai Medium,
  Scale       = 0.9
]

% Noto Sans Thai — clean modern sans-serif for structural elements
\newfontfamily\headingfont{Noto Sans Thai}[
  UprightFont = Noto Sans Thai Medium,
  BoldFont    = Noto Sans Thai Bold,
  Scale       = 0.9
]

% --- Heading font for structural elements ---
\setbeamerfont{title}{family=\headingfont, series=\bfseries, size=\huge}
\setbeamerfont{subtitle}{family=\headingfont, series=\mdseries}
\setbeamerfont{author}{family=\headingfont, series=\mdseries}

% --- Custom Title Page ---
\defbeamertemplate{title page}{custom}[1][]{%
  \centering
  \vspace*{2cm}
  % Title — in white
  {\usebeamerfont{title}\color{white}\inserttitle\par}
  \vspace{0.6cm}
  % Subtitle — in white
  {\usebeamerfont{subtitle}\color{white}\insertsubtitle\par}
  \vspace{2.5cm}
  % Author — blue filled rounded box, white text
  \begin{center}
    \begin{tcolorbox}[arc=3mm, colback=boxBlue, colframe=boxBlue!80,
                      width=0.42\textwidth, halign=center,
                      left=1.5em, right=1.5em, top=0.8em, bottom=0.8em,
                      boxrule=1.5pt, coltext=white]
      \usebeamerfont{author}\insertauthor
    \end{tcolorbox}
  \end{center}
}
\setbeamertemplate{title page}[custom]
\setbeamerfont{frametitle}{family=\headingfont, series=\bfseries}
\setbeamerfont{section in toc}{family=\headingfont}
\setbeamerfont{page number in head/foot}{family=\headingfont}
\setbeamerfont{section title}{family=\headingfont, series=\bfseries}
\setbeamerfont{subsection title}{family=\headingfont}

% Monospace font (for code/verbatim)
\setmonofont{Latin Modern Mono}[Scale=0.9]

% ---------- Math ----------
\usepackage{amsmath, amssymb}

% ---------- Graphics ----------
\usepackage{graphicx}
\usepackage{tikz}

% ---------- String parsing (for section headers) ----------
\usepackage{xstring}

% ---------- Colored boxes ----------
\usepackage[skins, breakable, xparse]{tcolorbox}
% Note: we do NOT load the 'most' option — it predefines example/solution environments
% that would conflict with our custom boxes.

% --- Color definitions ---
\definecolor{boxBlue}{RGB}{41, 128, 185}
\definecolor{boxGreen}{RGB}{39, 174, 96}
\definecolor{boxOrange}{RGB}{230, 126, 34}
\definecolor{boxPurple}{RGB}{142, 68, 173}
\definecolor{boxBlueBg}{RGB}{235, 245, 255}
\definecolor{boxGreenBg}{RGB}{235, 255, 240}
\definecolor{boxOrangeBg}{RGB}{255, 245, 235}
\definecolor{boxPurpleBg}{RGB}{248, 240, 255}

% --- Explanation box (blue) — for definitions, concepts, notes ---
\newtcolorbox{explanation}[1][]{
  enhanced,
  breakable,
  arc         = 2mm,
  colback     = boxBlueBg,
  colframe    = boxBlue!50,
  title       = #1,
  fonttitle   = \bfseries\color{white},
  coltitle    = white,
  colbacktitle= boxBlue,
  attach boxed title to top left = {yshift=-2mm, xshift=2mm},
  boxed title style = {arc=2mm, size=small},
  before skip = 8pt,
  after skip  = 8pt
}

% --- Example box (green) — for worked examples ---
\newtcolorbox{exambox}[1][]{
  enhanced,
  breakable,
  arc         = 2mm,
  colback     = boxGreenBg,
  colframe    = boxGreen!50,
  title       = #1,
  fonttitle   = \bfseries\color{white},
  coltitle    = white,
  colbacktitle= boxGreen,
  attach boxed title to top left = {yshift=-2mm, xshift=2mm},
  boxed title style = {arc=2mm, size=small},
  before skip = 8pt,
  after skip  = 8pt
}

% --- Solution box (orange) — for solutions and answers ---
\newtcolorbox{solbox}[1][]{
  enhanced,
  breakable,
  arc         = 2mm,
  colback     = boxOrangeBg,
  colframe    = boxOrange!50,
  title       = #1,
  fonttitle   = \bfseries\color{white},
  coltitle    = white,
  colbacktitle= boxOrange,
  attach boxed title to top left = {yshift=-2mm, xshift=2mm},
  boxed title style = {arc=2mm, size=small},
  before skip = 8pt,
  after skip  = 8pt
}

% --- Intuition box (purple) — for plain-language explanations, analogies ---
\newtcolorbox{intuition}[1][]{
  enhanced,
  breakable,
  arc         = 2mm,
  colback     = boxPurpleBg,
  colframe    = boxPurple!50,
  title       = #1,
  fonttitle   = \bfseries\color{white},
  coltitle    = white,
  colbacktitle= boxPurple,
  attach boxed title to top left = {yshift=-2mm, xshift=2mm},
  boxed title style = {arc=2mm, size=small},
  before skip = 8pt,
  after skip  = 8pt
}

% ---------- Beamer theme ----------
\usetheme{Madrid}
\usecolortheme{default}

% --- Custom beamer colors (blue accent matching explanation boxes) ---
\setbeamercolor{structure}{fg=boxBlue}
\setbeamercolor{palette primary}{fg=white, bg=boxBlue}
\setbeamercolor{palette secondary}{fg=white, bg=boxBlue!80}
\setbeamercolor{palette tertiary}{fg=white, bg=boxBlue!60}
\setbeamercolor{block title}{fg=white, bg=boxBlue}
\setbeamercolor{block body}{fg=black, bg=boxBlueBg}

% Remove navigation symbols for clean look
\setbeamertemplate{navigation symbols}{}

% Note: Frame title prefix removed — \addtobeamertemplate conflicts with beamer internals

% --- Outline (Table of Contents) styling ---
\setbeamerfont{section in toc}{family=\headingfont, series=\mdseries}
\setbeamerfont{subsection in toc}{family=\headingfont, series=\mdseries}
\setbeamertemplate{section in toc}{%
  \leavevmode\leftskip=1.5em
  \textcolor{boxBlue}{\textbf{\large\inserttocsectionnumber.}}%
  \quad\inserttocsection\par
  \vspace{0.2em}
}
\setbeamertemplate{subsection in toc}{%
  \leavevmode\leftskip=4em
  \textcolor{boxBlue!60}{\small$\bullet$}%
  \quad\inserttocsubsection\par
  \vspace{0.1em}
}

% Section header slides — Thai in white, English in smaller light blue below
\AtBeginSection{
  \begin{frame}[plain]{}
    \vspace*{2cm}
    \begin{beamercolorbox}[sep=12pt, center, rounded=true, shadow=true]{palette primary}
      \IfSubStr{\secname}{(}{%
        \StrBefore{\secname}{(}[\sectionthai]%
        \StrBetween[1,1]{\secname}{(}{)}[\sectioneng]%
        \begin{tabular}{@{}c@{}}
          {\usebeamerfont{title}\sectionthai} \\[0.2em]
          {\usebeamerfont{subtitle}\color{boxBlue!60}\sectioneng}
        \end{tabular}%
      }{%
        {\usebeamerfont{title}\secname\par}%
      }
    \end{beamercolorbox}
    \vfill
  \end{frame}
}

% Subsection header slides — smaller, centered
\AtBeginSubsection{
  \begin{frame}[plain]{}
    \vfill
    \centering
    {\usebeamerfont{section title}\insertsubsectionhead\par}
    \vfill
  \end{frame}
}

% Minimal footer: page number only, right-aligned
\setbeamertemplate{footline}{
  \hfill\usebeamerfont{page number in head/foot}\insertframenumber{} / \inserttotalframenumber\hspace*{1em}\vspace*{0.5ex}
}

% ---------- Spacing & readability ----------
\linespread{1.2}

% ---------- Useful commands ----------

% Highlight important text in blue
\newcommand{\highlight}[1]{\textcolor{boxBlue}{\textbf{#1}}}

% Math set notation helpers
\newcommand{\R}{\mathbb{R}}
\newcommand{\N}{\mathbb{N}}
\newcommand{\Z}{\mathbb{Z}}
\newcommand{\Q}{\mathbb{Q}}

% ============================================================================
%  End of preamble
% ============================================================================


% --- Document info ---
\title[ตรรกศาสตร์]{ตรรกศาสตร์พื้นฐาน}
\subtitle{Basic Logic — ประพจน์, ตัวเชื่อม, ตัวบ่งปริมาณ, การอ้างเหตุผล}
\author[None W. (Yuan)]{None W. (Yuan)}
\date{}

\begin{document}

% ===========================================================================
%  TITLE SLIDE
% ===========================================================================
\begin{frame}[plain]
  \titlepage
\end{frame}

% ===========================================================================
%  OUTLINE
% ===========================================================================
\begin{frame}{สารบัญ (Outline)}
  \tableofcontents
\end{frame}

% ===========================================================================
%  ก่อนเรียน (Prerequisites)
% ===========================================================================
\begin{frame}{ก่อนเรียน (Prerequisites)}
  \begin{explanation}[สิ่งที่ควรรู้ก่อนเรียน]
    \begin{itemize}
      \item \textbf{เซต} — การดำเนินการของเซต ($\cup$, $\cap$, $'$) ซึ่งเชื่อมโยงกับตัวเชื่อมทางตรรกศาสตร์
    \end{itemize}
  \end{explanation}

  \begin{intuition}[ตรรกศาสตร์คืออะไร?]
    ตรรกศาสตร์คือการศึกษา \highlight{การให้เหตุผล} — อะไรจริง อะไรเท็จ
    และการสรุปผลจากสิ่งที่รู้

    \medskip
    เป็นพื้นฐานของการเขียนโปรแกรม (if-else, conditions), การพิสูจน์ทางคณิตศาสตร์,
    และการแก้ปัญหาการจับโกหก!
  \end{intuition}
\end{frame}

% ===========================================================================
%  SECTION 1: ประพจน์และตัวเชื่อม
% ===========================================================================
\section{ประพจน์และตัวเชื่อม (Propositions \& Connectives)}

\subsection{ประพจน์}

\begin{frame}{ประพจน์คืออะไร?}
  \begin{explanation}[นิยาม: ประพจน์ (Proposition)]
    \textbf{ประพจน์} คือประโยคบอกเล่าหรือปฏิเสธที่มีค่าความจริงเป็น \highlight{จริง (T)} หรือ \highlight{เท็จ (F)} อย่างใดอย่างหนึ่งเท่านั้น
  \end{explanation}

  \begin{columns}[T]
    \column{0.5\textwidth}
      \textbf{เป็นประพจน์}
      \begin{itemize}
        \item $1 + 1 = 2$ \quad (จริง)
        \item กรุงเทพเป็นเมืองหลวงของไทย \quad (จริง)
        \item $2 > 5$ \quad (เท็จ)
      \end{itemize}

    \column{0.5\textwidth}
      \textbf{ไม่เป็นประพจน์}
      \begin{itemize}
        \item ``ไปไหนมา?'' (คำถาม)
        \item ``กรุณาปิดประตู'' (คำสั่ง)
        \item ``$x + 1 = 3$'' (มีตัวแปรอิสระ — ไม่ทราบจริง/เท็จ)
      \end{itemize}
  \end{columns}
\end{frame}

\subsection{ตัวเชื่อม (Logical Connectives)}

\begin{frame}{นิเสธ (Negation) — NOT}
  \begin{explanation}[นิยาม]
    \textbf{นิเสธ} ของ $p$ ($\lnot p$ หรือ $\sim p$) มีค่าความจริง \highlight{ตรงข้าม}กับ $p$

    \medskip
    \begin{center}
      \begin{tabular}{c|c}
        $p$ & $\lnot p$ \\
        \hline
        T & F \\
        F & T
      \end{tabular}
    \end{center}
  \end{explanation}

  \begin{exambox}[ตัวอย่าง]
    $p$: ``วันนี้ฝนตก''

    $\lnot p$: ``วันนี้ฝนไม่ตก''

    \medskip
    ถ้า $p$ เป็นจริง (ฝนตกจริง) $\rightarrow$ $\lnot p$ เป็นเท็จ
  \end{exambox}
\end{frame}

\begin{frame}{ตัวเชื่อม ``และ'' (Conjunction) — AND}
  \begin{explanation}[นิยาม]
    $p \land q$ (อ่านว่า ``$p$ และ $q$'') เป็นจริง \highlight{ก็ต่อเมื่อทั้ง $p$ และ $q$ เป็นจริงทั้งคู่}

    \medskip
    \begin{center}
      \begin{tabular}{cc|c}
        $p$ & $q$ & $p \land q$ \\
        \hline
        T & T & T \\
        T & F & F \\
        F & T & F \\
        F & F & F
      \end{tabular}
    \end{center}
  \end{explanation}

  \begin{exambox}[ตัวอย่าง]
    $p$: ``ฉันมีเงิน'' (T), \quad $q$: ``ร้านเปิด'' (F)

    $p \land q$: ``ฉันมีเงิน \textbf{และ} ร้านเปิด'' $\rightarrow$ F (เพราะ $q$ เป็น F)
  \end{exambox}
\end{frame}

\begin{frame}{ตัวเชื่อม ``หรือ'' (Disjunction) — OR}
  \begin{explanation}[นิยาม]
    $p \lor q$ (อ่านว่า ``$p$ หรือ $q$'') เป็นเท็จ \highlight{ก็ต่อเมื่อทั้ง $p$ และ $q$ เป็นเท็จทั้งคู่}

    \medskip
    \begin{center}
      \begin{tabular}{cc|c}
        $p$ & $q$ & $p \lor q$ \\
        \hline
        T & T & T \\
        T & F & T \\
        F & T & T \\
        F & F & F
      \end{tabular}
    \end{center}
  \end{explanation}

  \begin{intuition}[ข้อสังเกต]
    ในตรรกศาสตร์ ``หรือ'' หมายถึง \highlight{หรืออย่างน้อยหนึ่งอย่าง}
    (ไม่ใช่ ``หรืออย่างใดอย่างหนึ่งเท่านั้น'' เหมือนในชีวิตประจำวัน)
  \end{intuition}
\end{frame}

\begin{frame}{ตัวเชื่อม ``ถ้า...แล้ว...'' (Implication) — IF-THEN}
  \begin{explanation}[นิยาม]
    $p \to q$ (อ่านว่า ``ถ้า $p$ แล้ว $q$'') เป็นเท็จ \highlight{ก็ต่อเมื่อ $p$ จริง แต่ $q$ เท็จ}

    \medskip
    \begin{center}
      \begin{tabular}{cc|c}
        $p$ & $q$ & $p \to q$ \\
        \hline
        T & T & T \\
        T & F & F \\
        F & T & T \\
        F & F & T
      \end{tabular}
    \end{center}
  \end{explanation}

  \begin{intuition}[จำง่าย ๆ]
    ``ถ้าสัญญาแล้วต้องทำ'' — ผิดสัญญาก็ต่อเมื่อ \highlight{สัญญา (T) แต่ไม่ทำ (F)} เท่านั้น

    ถ้าไม่ได้สัญญา ($p$ = F) แล้วทำ ($q$ = T) ก็ไม่ผิดสัญญา
  \end{intuition}
\end{frame}

\begin{frame}{ตัวเชื่อม ``ก็ต่อเมื่อ'' (Biconditional) — IFF}
  \begin{explanation}[นิยาม]
    $p \leftrightarrow q$ (อ่านว่า ``$p$ ก็ต่อเมื่อ $q$'') เป็นจริง \highlight{เมื่อ $p$ และ $q$ มีค่าความจริงเหมือนกัน}

    \medskip
    \begin{center}
      \begin{tabular}{cc|c}
        $p$ & $q$ & $p \leftrightarrow q$ \\
        \hline
        T & T & T \\
        T & F & F \\
        F & T & F \\
        F & F & T
      \end{tabular}
    \end{center}
  \end{explanation}

  \begin{intuition}[ความหมาย]
    $p \leftrightarrow q$ เทียบเท่ากับ $(p \to q) \land (q \to p)$

    ``$p$ เป็นจริง \textbf{ก็ต่อเมื่อ} $q$ เป็นจริง'' = ทั้งคู่จริงพร้อมกัน หรือเท็จพร้อมกัน
  \end{intuition}
\end{frame}

\begin{frame}{สรุปตารางค่าความจริง}
  \begin{center}
    \begin{tabular}{cc|c|c|c|c|c}
      $p$ & $q$ & $\lnot p$ & $p \land q$ & $p \lor q$ & $p \to q$ & $p \leftrightarrow q$ \\
      \hline
      T & T & F & T & T & T & T \\
      T & F & F & F & T & F & F \\
      F & T & T & F & T & T & F \\
      F & F & T & F & F & T & T
    \end{tabular}
  \end{center}

  \begin{intuition}[วิธีจำ]
    \begin{itemize}
      \item $\land$ (และ): เป็นจริงแค่กรณีเดียว — T,T
      \item $\lor$ (หรือ): เป็นเท็จแค่กรณีเดียว — F,F
      \item $\to$ (ถ้า...แล้ว): เป็นเท็จแค่กรณีเดียว — T,F
      \item $\leftrightarrow$ (ก็ต่อเมื่อ): เหมือนกัน = จริง, ต่างกัน = เท็จ
    \end{itemize}
  \end{intuition}
\end{frame}

% ===========================================================================
%  SECTION 2: การสมมูลกันของประพจน์ (Logical Equivalence)
% ===========================================================================
\section{การสมมูลกันของประพจน์ (Logical Equivalence)}

\begin{frame}{ประพจน์ที่สมมูลกัน}
  \begin{explanation}[นิยาม]
    ประพจน์สองรูป \textbf{สมมูลกัน} ($P \equiv Q$) ก็ต่อเมื่อ
    \highlight{มีค่าความจริงเหมือนกันในทุกกรณี}
  \end{explanation}

  \begin{exambox}[ตัวอย่าง: Contrapositive]
    $p \to q \;\equiv\; \lnot q \to \lnot p$

    \medskip
    \begin{center}
      \begin{tabular}{cc|c|c}
        $p$ & $q$ & $p \to q$ & $\lnot q \to \lnot p$ \\
        \hline
        T & T & T & T \\
        T & F & F & F \\
        F & T & T & T \\
        F & F & T & T
      \end{tabular}
    \end{center}

    ค่าความจริงตรงกันทุกกรณี $\rightarrow$ สมมูลกัน $\checkmark$
  \end{exambox}
\end{frame}

\begin{frame}{กฎที่สำคัญ}
  \begin{explanation}[กฎเดอมอร์แกน (De Morgan's Laws)]
    \begin{enumerate}
      \item $\lnot (p \land q) \equiv \lnot p \lor \lnot q$
      \item $\lnot (p \lor q) \equiv \lnot p \land \lnot q$
    \end{enumerate}
  \end{explanation}

  \begin{explanation}[กฎอื่น ๆ]
    \begin{itemize}
      \item \textbf{Double Negation}: $\lnot (\lnot p) \equiv p$
      \item \textbf{Implication}: $p \to q \equiv \lnot p \lor q$
      \item \textbf{Contrapositive}: $p \to q \equiv \lnot q \to \lnot p$
    \end{itemize}
  \end{explanation}
\end{frame}

\begin{frame}{ตรรกศาสตร์กับเซต}
  \begin{intuition}[เชื่อมโยงระหว่างตรรกศาสตร์กับเซต]
    \[
    \begin{array}{c|c}
      \textbf{ตรรกศาสตร์} & \textbf{เซต} \\
      \hline
      \lnot \text{ (นิเสธ)} & ' \text{ (คอมพลีเมนต์)} \\
      \land \text{ (และ)} & \cap \text{ (อินเตอร์เซกชัน)} \\
      \lor \text{ (หรือ)} & \cup \text{ (ยูเนียน)} \\
      T & U \text{ (เอกภพสัมพัทธ์)} \\
      F & \emptyset \text{ (เซตว่าง)}
    \end{array}
    \]

    \medskip
    กฎเดอมอร์แกนใช้ได้เหมือนกันทั้งสองระบบ!
  \end{intuition}
\end{frame}

% ===========================================================================
%  SECTION 3: ตัวบ่งปริมาณ (Quantifiers)
% ===========================================================================
\section{ตัวบ่งปริมาณ (Quantifiers)}

\begin{frame}{ตัวบ่งปริมาณคืออะไร?}
  \begin{explanation}[นิยาม]
    \textbf{ตัวบ่งปริมาณ} ใช้กับประโยคเปิด (ประโยคที่มีตัวแปร) เพื่อบอกขอบเขตของตัวแปร

    \medskip
    \begin{itemize}
      \item $\forall x\; P(x)$ = \highlight{``สำหรับทุก $x$, $P(x)$ เป็นจริง''}
      \item $\exists x\; P(x)$ = \highlight{``มี $x$ อย่างน้อย 1 ตัว ที่ $P(x)$ เป็นจริง''}
    \end{itemize}
  \end{explanation}

  \begin{exambox}[ตัวอย่าง]
    \begin{itemize}
      \item $\forall x \in \R:\; x^2 \geq 0$ \quad (จริง — กำลังสองของจำนวนจริง $\geq 0$ เสมอ)
      \item $\exists x \in \Z:\; x^2 = 4$ \quad (จริง — $x = 2$ หรือ $x = -2$)
      \item $\forall x \in \R:\; x^2 = 4$ \quad (เท็จ — ไม่ใช่ทุกจำนวนจริงที่ยกกำลังสองแล้วได้ 4)
    \end{itemize}
  \end{exambox}
\end{frame}

\begin{frame}{นิเสธของตัวบ่งปริมาณ}
  \begin{explanation}[กฎการกระจายนิเสธ]
    \begin{enumerate}
      \item $\lnot (\forall x\; P(x)) \;\equiv\; \exists x\; \lnot P(x)$
      \item $\lnot (\exists x\; P(x)) \;\equiv\; \forall x\; \lnot P(x)$
    \end{enumerate}
  \end{explanation}

  \begin{exambox}[ตัวอย่าง]
    นิเสธของ ``นกทุกตัวบินได้'' คืออะไร?

    \medskip
    ประโยคเดิม: $\forall x\; (x \text{ บินได้})$ \quad ($x$ = นก)

    นิเสธ: $\lnot \forall x\; (x \text{ บินได้}) \equiv \exists x\; \lnot(x \text{ บินได้})$

    \medskip
    ``\textbf{มี}นกอย่างน้อย 1 ตัวที่\textbf{บินไม่ได้}'' \quad (เช่น นกกระจอกเทศ)
  \end{exambox}
\end{frame}

% ===========================================================================
%  SECTION 4: การอ้างเหตุผล (Rules of Inference)
% ===========================================================================
\section{การอ้างเหตุผล (Rules of Inference)}

\begin{frame}{Modus Ponens และ Modus Tollens}
  \begin{explanation}[Modus Ponens (การยืนยัน)]
    \[
      \begin{array}{c}
        p \to q \\
        p \\
        \hline
        \therefore\; q
      \end{array}
    \]

    ``ถ้า $p$ แล้ว $q$'' และ $p$ เป็นจริง $\rightarrow$ สรุปว่า $q$ เป็นจริง
  \end{explanation}

  \begin{explanation}[Modus Tollens (การปฏิเสธ)]
    \[
      \begin{array}{c}
        p \to q \\
        \lnot q \\
        \hline
        \therefore\; \lnot p
      \end{array}
    \]

    ``ถ้า $p$ แล้ว $q$'' และ $q$ เป็นเท็จ $\rightarrow$ สรุปว่า $p$ เป็นเท็จ
  \end{explanation}
\end{frame}

\begin{frame}{ตัวอย่าง: Modus Ponens \& Modus Tollens}
  \begin{exambox}[Modus Ponens]
    \textbf{เหตุ:}
    \begin{enumerate}
      \item ถ้าฝนตก แล้วถนนจะเปียก \quad ($p \to q$)
      \item ฝนตก \quad ($p$)
    \end{enumerate}
    \textbf{ผล:} ถนนเปียก \quad ($q$) \quad $\checkmark$
  \end{exambox}

  \begin{exambox}[Modus Tollens]
    \textbf{เหตุ:}
    \begin{enumerate}
      \item ถ้าเป็นนกเพนกวิน แล้วบินไม่ได้ \quad ($p \to q$)
      \item เจ้านี่บินได้ \quad ($\lnot q$)
    \end{enumerate}
    \textbf{ผล:} เจ้านี่ไม่ใช่นกเพนกวิน \quad ($\lnot p$) \quad $\checkmark$
  \end{exambox}
\end{frame}

\begin{frame}{ระวัง! การอ้างเหตุผลที่ผิด}
  \begin{exambox}[ตัวอย่างการอ้างเหตุผลที่ \textbf{ไม่} สมเหตุสมผล]
    \textbf{Converse Error (ยืนยันผล):}
    \begin{enumerate}
      \item ถ้าฝนตก แล้วถนนเปียก \quad ($p \to q$)
      \item ถนนเปียก \quad ($q$)
    \end{enumerate}
    สรุปว่า ฝนตก \quad \textbf{ผิด!}

    \medskip
    (ถนนอาจเปียกเพราะรถน้ำผ่าน ไม่จำเป็นต้องฝนตก)

    \medskip
    \textbf{Inverse Error (ปฏิเสธเหตุ):}
    \begin{enumerate}
      \item ถ้าฝนตก แล้วถนนเปียก \quad ($p \to q$)
      \item ฝนไม่ตก \quad ($\lnot p$)
    \end{enumerate}
    สรุปว่า ถนนไม่เปียก \quad \textbf{ผิด!}
  \end{exambox}
\end{frame}

\begin{frame}{ตัวอย่าง: ข้อสอบจริง}
  \begin{exambox}[ข้อสอบปี 67 ข้อ 17]
    ถ้าวันนี้ร้อน และอบอ้าว แล้วฝนตก ถ้าวันนี้ฝนไม่ตก ข้อใดถูกต้อง
  \end{exambox}

  \begin{solbox}[วิธีทำ]
    ให้ $p = $ วันนี้ร้อน และอบอ้าว $q = $ วันนี้ฝนตก \newline
    จากข้อนี้เหตุ (สิ่งที่เกิดขึ้นแล้ว) คือ 
    \begin{enumerate}
      \item $p \rightarrow q$
      \item $\lnot q$
    \end{enumerate}
    จาก Modus Tollens ผลคือ $\lnot p$ ซึ่งเมื่อใช้กฎของ De Morgan จะได้เป็น \textbf{วันนี้ไม่ร้อน หรือ ไม่อบอ้าว}
  \end{solbox}
\end{frame}

\begin{frame}{วิธีทำโจทย์การอ้างเหตุผล}
  \begin{explanation}[วิธีทำโจทย์การอ้างเหตุผล]
    \begin{enumerate}
      \item สร้างประพจน์ใหม่ ดังนี้ ถ้า เหตุทั้งหมดมาต่อกันด้วย \textbf{และ $\wedge$} แล้ว ผล \newline
      (เหตุ$_1$ $\wedge$ เหตุ$_2$ $\wedge \cdots  \wedge $ เหตุ$_n$) $\rightarrow$ ผล
      \item พิสูจน์ว่าประพจน์นี้เป็น \textbf{สัจนิรันดร์ (Tautology)} คือมีค่าความจริงเป็นจริงเสมอ
      \begin{itemize}
        \item ถ้าเป็น \textbf{สัจนิรันดร์} สมเหตุสมผล 
        \item ถ้า \textbf{ไม่เป็นสัจนิรันดร์} ไม่สมเหตุสมผล
      \end{itemize}
      \item วิธีการพิสูจน์สัจนิรันดร์คือไม่ว่าประพจน์ย่อยจะมีค่าความจริงแบบใด ประพจน์ใหม่จะมีค่าความจริงเป็นจริงเสมอ
    \end{enumerate}
  \end{explanation}
\end{frame}

\begin{frame}{วิธีทำโจทย์การอ้างเหตุผล}
  \begin{intuition}[วิธีทำโจทย์การอ้างเหตุผลแบบเร็ว]
    \begin{enumerate}
      \item ให้เหตุทั้งหมดมีค่าความจริงเป็น \textbf{จริง}
      \item จากนั้นหาค่าความจริงของประพจน์ย่อยแต่ละประพจน์เท่าที่หาได้
      \item จากนั้นตรวจสอบว่าค่าความจริงของผลจากค่าความจริงของประพจน์ย่อยที่ได้จากเหตุ 
      \begin{itemize}
        \item ถ้าเป็น \textbf{จริง} สมเหตุสมผล 
        \item ถ้าเป็น \textbf{เท็จ} ไม่สมเหตุสมผล
        \item ถ้าไม่ทราบค่าความจริง ไม่สามารถสรุปได้ด้วยวิธีนี้ \textbf{*}
      \end{itemize}
    \end{enumerate}
  \end{intuition}

  \begin{exambox}[*]
    บางกรณีไม่ต้องทราบค่าความจริงทั้งหมดก็สรุปได้ เช่น ถ้าเหตุเป็น ถ้าฝนตกแล้วถนนเปียก แล้วผลเป็น ถ้าฝนตกแล้วถนนเปียก ต้องใช้วิธีการหาสัจนิรันดร์
  \end{exambox}
\end{frame}

\begin{frame}{ตัวอย่าง: การอ้างเหตุผลหลายเหตุ}
  \begin{exambox}[โจทย์]
    พิจารณาการอ้างเหตุผลต่อไปนี้สมเหตุสมผลหรือไม่

    \textbf{เหตุ:}
    \begin{enumerate}
      \item ถ้าฝนตก แล้วถนนเปียก \quad ($p \to q$)
      \item ถ้าถนนเปียก แล้วการจราจรติด \quad ($q \to r$)
      \item ฝนตก \quad ($p$)
    \end{enumerate}
    \textbf{ผล:} การจราจรติด \quad ($r$)
  \end{exambox}
\end{frame}

\begin{frame}{ตัวอย่าง: การอ้างเหตุผลหลายเหตุ (ต่อ)}
  \begin{solbox}[วิธีทำ — แบบที่ 1: ใช้ Modus Ponens]
    \begin{align*}
      &(1)\; p \to q    &&\text{(เหตุข้อ 1)} \\
      &(2)\; q \to r    &&\text{(เหตุข้อ 2)} \\
      &(3)\; p          &&\text{(เหตุข้อ 3)} \\
      &(4)\; q          &&\text{(จาก (1) และ (3) ด้วย Modus Ponens)} \\
      &(5)\; r          &&\text{(จาก (2) และ (4) ด้วย Modus Ponens)}
    \end{align*}

    ได้ $r$ ตรงกับผล $\rightarrow$ \textbf{สมเหตุสมผล} $\checkmark$
  \end{solbox}
\end{frame}

\begin{frame}{ตัวอย่าง: การอ้างเหตุผลหลายเหตุ (ต่อ)}
  \begin{solbox}[วิธีทำ — แบบที่ 2: สัจนิรันดร์]
    สร้างประพจน์: $[(p \to q) \land (q \to r) \land p] \to r$

    ตรวจสอบโดยให้เหตุเป็นจริงทั้งหมด:
    \begin{itemize}
      \item $p$ เป็นจริง
      \item $p \to q$ เป็นจริง และ $p$ เป็นจริง $\rightarrow$ $q$ ต้องเป็นจริง
      \item $q \to r$ เป็นจริง และ $q$ เป็นจริง $\rightarrow$ $r$ ต้องเป็นจริง
    \end{itemize}
    เมื่อเหตุเป็นจริง ผล ($r$) ก็เป็นจริงเสมอ $\rightarrow$ \textbf{สมเหตุสมผล} $\checkmark$
  \end{solbox}

  \begin{intuition}[Hypothetical Syllogism]
    จาก $p \to q$ และ $q \to r$ สามารถสรุป $p \to r$ ได้ทันที
    (กฎการถ่ายทอด — Transitivity of Implication)
  \end{intuition}
\end{frame}

% ===========================================================================
%  SECTION 5: ปัญหาการจับโกหก (Liar Puzzles)
% ===========================================================================
\section{ปัญหาการจับโกหก (Liar Puzzles)}

\begin{frame}{แนวคิดการแก้ปัญหาจับโกหก}
  \begin{explanation}[หลักการ]
    ในปัญหาจับโกหก เรามีบุคคลที่:
    \begin{itemize}
      \item พูด\textbf{จริง}เสมอ (Truth-teller) — ทุกประโยคที่พูดมีค่าเป็น T
      \item พูด\textbf{เท็จ}เสมอ (Liar) — ทุกประโยคที่พูดมีค่าเป็น F
    \end{itemize}

    \medskip
    \textbf{วิธีแก้:} \highlight{สมมติให้แต่ละคนเป็นคนพูดจริง (หรือพูดเท็จ)}
    แล้วตรวจสอบความขัดแย้ง
  \end{explanation}

  \begin{exambox}[ตัวอย่างคลาสสิก]
    A พูดว่า ``ฉันพูดเท็จ''

    \medskip
    \begin{itemize}
      \item ถ้า A พูดจริง $\rightarrow$ ประโยค ``ฉันพูดเท็จ'' เป็นจริง $\rightarrow$ ขัดแย้ง!
      \item ถ้า A พูดเท็จ $\rightarrow$ ประโยค ``ฉันพูดเท็จ'' เป็นเท็จ $\rightarrow$ A พูดจริง $\rightarrow$ ขัดแย้ง!
    \end{itemize}

    นี่คือ \textbf{Paradox} — เป็นไปไม่ได้ทั้งสองกรณี!
  \end{exambox}
\end{frame}

\begin{frame}{ตัวอย่าง: จับโกหก}
  \begin{exambox}[โจทย์]
    มี 2 คน: A และ B

    A พูดว่า ``พวกเราทั้งคู่เป็นคนพูดเท็จ''

    จงระบุว่า A และ B เป็นคนประเภทใด?
  \end{exambox}

  \begin{solbox}[วิธีทำ]
    \textbf{กรณี 1: A พูดจริง}
    \begin{itemize}
      \item ``ทั้งคู่เป็นคนพูดเท็จ'' เป็นจริง $\rightarrow$ A พูดเท็จ \textbf{ขัดแย้ง!}
    \end{itemize}
    ดังนั้น A \textbf{ต้อง}เป็นคนพูดเท็จ
  \end{solbox}
\end{frame}

\begin{frame}{ตัวอย่าง: จับโกหก (ต่อ)}
  \begin{solbox}[วิธีทำ (ต่อ)]
    \textbf{กรณี 2: A พูดเท็จ}
    \begin{itemize}
      \item ``ทั้งคู่เป็นคนพูดเท็จ'' เป็นเท็จ
      \item ดังนั้นมีอย่างน้อย 1 คนที่พูดจริง
      \item A พูดเท็จแล้ว $\rightarrow$ B ต้องพูดจริง
    \end{itemize}

    \textbf{สรุป:} A เป็นคนพูดเท็จ, B เป็นคนพูดจริง
  \end{solbox}

  \begin{intuition}[เทคนิค]
    วิธีแก้ปัญหาจับโกหก: \highlight{สมมติทีละกรณี} แล้วหาข้อขัดแย้ง
    กรณีไหนไม่ขัดแย้ง = คำตอบที่ถูกต้อง
  \end{intuition}
\end{frame}

% ===========================================================================
%  SUMMARY
% ===========================================================================
\section{สรุป}

\begin{frame}{สรุปเนื้อหา}
  \begin{explanation}[สิ่งที่ได้เรียนรู้]
    \begin{enumerate}
      \item \textbf{ประพจน์}: ประโยคที่มีค่าความจริง T/F
      \item \textbf{ตัวเชื่อม}: $\lnot$ (ไม่), $\land$ (และ), $\lor$ (หรือ), $\to$ (ถ้า...แล้ว), $\leftrightarrow$ (ก็ต่อเมื่อ)
      \item \textbf{การสมมูล}: De Morgan, Contrapositive, Implication
      \item \textbf{ตัวบ่งปริมาณ}: $\forall$ (สำหรับทุก), $\exists$ (มีอย่างน้อยหนึ่ง)
      \item \textbf{การอ้างเหตุผล}: Modus Ponens, Modus Tollens
      \item \textbf{ปัญหาจับโกหก}: สมมติ + ตรวจสอบความขัดแย้ง
    \end{enumerate}
  \end{explanation}

  \begin{center}
    \large เอกสารนี้เป็นส่วนหนึ่งของเนื้อหาเตรียมสอบ \textbf{สอวน. คอมพิวเตอร์}
  \end{center}
\end{frame}

\end{document}
